\chapter*{ABSTRACT}
\addcontentsline{toc}{chapter}{Abstract}

Many everyday conversations are attempts at persuasion: a fundraiser asking for a donation, or an online seller negotiating with a buyer over chat. These conversations are asymmetric. One person argues and makes offers, while the other weighs them and finally decides. For the people running such conversations at scale, two questions matter most: is the other person actually going to say yes, and which persuasion tactics are being used along the way? Existing text classifiers usually treat a dialogue as one undivided block of text, which ignores who is speaking, cuts long conversations short, and often leans on surface cues in the last few lines.

This project, \textit{Antiphony}, builds a framework that answers both questions. For commitment prediction, we manually re-labeled the \textit{Persuasion for Good} charity dialogues, because the dataset's recorded donation amounts often disagree with what participants actually said. We then designed the \textit{AntiphonyClassifier}, a model that reads the two speakers as separate streams and lets the decision-maker's side attend to the arguer's side, following how the decision builds up over the conversation. It was the strongest of the fourteen model versions we built, at about $0.86$ Macro-F1. For strategy recognition, we labeled every persuader turn with the tactics it uses under an $11$-category, $41$-strategy scheme, and built an ensemble of classifiers that gets close to the agreement limit of the labels themselves.

To test whether the approach carries over to commerce, we created a new 700-dialogue buyer--seller negotiation benchmark, including dialogues in which the buyer changes their mind before the end. On these, Antiphony holds its accuracy better than a model that only reads the final turns, which shows that it tracks the whole course of the decision rather than its last words.

The resulting models are small enough to run on a single commodity GPU. They can help fundraising and sales teams spot likely commitments early, prioritize promising leads, review which tactics work, and flag manipulative pressure in automated conversations.

\vspace{12pt}
\noindent\textbf{Keywords:} Persuasion Strategy, Conversational Intention, Asymmetric Dialogue, Antiphony Architecture, Multi-Label Classification, Adversarial Reversal.
\clearpage
